\section{Marco teórico}
Para comprender el funcionamiento del canal de ECG implementado en esta práctica, es necesario analizar los principales conceptos relacionados con el acondicionamiento de señales biomédicas: el amplificador de instrumentación AD620, el rechazo de modo común (CMRR) y su aplicación en la adquisición de un electrocardiograma.

\subsection{AD620}
El AD620 es un amplificador de instrumentación integrado cuyo diseño tiene un doble objetivo: la amplificación precisa de señales diferenciales de pequeña amplitud y la supresión de señales de modo común. 
Este amplificador es ampliamente utilizado en aplicaciones biomédicas debido a su alta impedancia de entrada, bajo nivel de ruido y capacidad para operar con una amplia gama de voltajes de alimentación.
Estas características hacen que el amplificador de instrumentación sea apropiado para aplicaciones como la adquisición de señales de ECG.

Una de las principales características del AD620 es su alta impedancia de entrada, especificada en el orden de los $10~\mathrm{G}\Omega$.
 Esto resulta particularmente importante en aplicaciones biomédicas, ya que los electrodos utilizados para adquirir señales fisiológicas presentan una impedancia propia.
 Si la impedancia de entrada del sistema de medición fuera comparable con la impedancia electrodo-piel, podría producirse un divisor resistivo que atenuara la señal antes de que esta fuera amplificada.
 Una impedancia de entrada elevada minimiza este efecto y permite obtener una medición más acertada de la señal biológica.
 
Otra ventaja del AD620 es que su ganancia puede configurarse mediante una única resistencia externa $R_G$, conectada entre los terminales de ganancia del integrado.
La relación entre la ganancia y dicha resistencia está dada por:
\begin{equation}
G = 1 + \frac{49{,}4\,\mathrm{k}\Omega}{R_G}
\end{equation}
 Por lo tanto, al disminuir el valor de $R_G$, aumenta la ganancia del amplificador, lo que permite adaptar la amplificación a la amplitud de la señal que se desea medir.
 En el caso del ECG, resulta necesario seleccionar cuidadosamente la ganancia para que la señal pueda ser observada y procesada sin amplificar excesivamente el ruido ni reducir significativamente el ancho de banda.

\subsection{CMRR}
En una medición es necesario distinguir entre las señales diferenciales y las señales de modo común. La señal diferencial es aquella que buscamos medir y corresponde a la diferencia de potencial entre los dos terminales de entrada del amplificador según la siguiente expresión:
\begin{equation} 
V_D = V_1 - V_2 
\end{equation}
Por el otro lado, las señales de modo común son aquellas que aparecen simultáneamente en ambos
 terminales de entrada con la misma amplitud, fase y frecuencia y se escriben como: 
\begin{equation} 
V_{CM} = \frac{V_1 + V_2}{2}
\end{equation}

La capacidad de un amplificador para rechazar las señales de modo común se caracteriza mediante el CMRR (\textit{Common Mode Rejection Ratio}), definido como la relación entre la ganancia diferencial y la ganancia en modo común: 
\begin{equation} 
CMRR = \frac{G_D}{G_{CM}}
\end{equation} 
También puede expresarse en decibeles como: 
\begin{equation}
CMRR_{\mathrm{dB}} = 20\log_{10}\left(\frac{G_D}{G_{CM}}\right)
\end{equation} 
Un CMRR elevado indica que el amplificador amplifica principalmente la diferencia de potencial entre sus entradas, mientras que atenúa las señales presentes simultáneamente en ambas. 
Esta característica resulta especialmente importante en aplicaciones biomédicas, donde pueden existir interferencias externas que se acoplan de manera similar a los dos electrodos. 

\subsection{Aplicación en ECG}
El electrocardiograma es el registro de la actividad eléctrica del corazón obtenido mediante electrodos colocados sobre la superficie corporal.
 La señal medida presenta una amplitud pequeña, del orden de los milivolts, por lo que su adquisición requiere una etapa de acondicionamiento capaz de amplificarla sin introducir distorsiones significativas y de reducir las interferencias presentes durante la medición.
 Para obtener el ECG se utilizan electrodos superficiales ubicados en brazo derecho, brazo izquierdo y pierna izquierda, que permiten visualizar las derivaciones DI, DII y DIII.
 
Por lo tanto, debido a la pequeña amplitud de la señal cardíaca, el AD620 se utiliza como etapa de preamplificación diferencial, realzando la señal de interés y rechazando las interferencias de modo común presentes en el entorno pero sin aumentar de manera tan significativa la ganacia.
 La ganancia se aumenta posteriormente mediante un amplificador operacional en configuración no inversora, que permite obtener la señal final con la amplitud adecuada para su visualización y procesamiento.
 Luego de la preamplificación es necesario incorporar etapas de filtrado que permitan conservar las componentes relevantes de la señal cardíaca y atenuar aquellas que no forman parte de la información de interés.
 En particular, un filtro pasa altos permite reducir la componente continua y las variaciones lentas de la línea de base, mientras que un filtro pasa bajos permite atenuar componentes de mayor frecuencia y parte del ruido presente en la medición.

En resumen, el canal de ECG es una cadena de acondicionamiento: los electrodos captan la señal, el AD620 amplifica la componente diferencial y rechaza el modo común, y las etapas posteriores de amplificación y filtrado la adecuan para su análisis. 